\documentclass[11pt,a4paper]{article}
\usepackage[margin=0.85in,top=1in,bottom=1in]{geometry}
\usepackage{amsmath,amssymb,amsfonts}
\usepackage{graphicx}
\usepackage{float}
\usepackage{enumitem}
\usepackage{microtype}
\usepackage{caption}
\usepackage[hidelinks]{hyperref}
\usepackage[most,breakable]{tcolorbox}
\tcbuselibrary{breakable,skins}
\usepackage{fontspec}
\setmainfont{DejaVu Serif}
\setsansfont{DejaVu Sans}
\setmonofont{DejaVu Sans Mono}
\captionsetup{font=small,labelfont=bf,skip=4pt}
\setlist{leftmargin=1.5em,topsep=2pt}

%% Breakable problem card — prevents content cutoff across pages
\newtcolorbox{solcard}[3]{
  breakable, enhanced,
  colback=white, colframe=gray!50, arc=2pt,
  title={\textbf{#1}\hfill\textsf{\small #3\,pts}},
  fonttitle=\normalsize\sffamily\bfseries,
  before upper={\small\textit{#2}\par\smallskip\hrule height 0.4pt\smallskip},
  top=4pt, bottom=6pt, left=7pt, right=7pt,
  parbox=false,
}

\title{\textbf{Long line 2.0 --- Solution}}
\author{\normalsize Y25 (2025) — English translation}
\date{}
\begin{document}\maketitle\vspace{-0.5em}
\begin{solcard}{A1}{Express the voltage across resistor $U_R(t)$, and the current through it $I_R(t)$ in terms of $U_l(x, t)$, $U_r(x, t)$, $I_l(x, t)$, $I_r(x, t)$ and $\mathcal E(t)$. Obtain an expression connecting $U_l(0, t)$ and $U_r(0, t)$. The expression may contain $R$, $Z$, $\mathcal E(t)$.}{0.40}

The voltage drop across the resistor is $U_R(t)=\mathcal E(t)-(U_l(0,t)+U_r(0,t))$ and the current is $I_R(t)=\dfrac{U_R(t)}{R}=I_l(0,t) + I_r(0,t)=\dfrac{1}{Z}\left(U_r(0,t)-U_l(0,t)\right)$. From here we get: \[\dfrac{\mathcal E(t)}{R}-\dfrac{U_l(0,t)}{R}-\dfrac{U_r(0,t)}{R}=\dfrac{U_r(0,t)}{Z}-\dfrac{U_l(0,t)}{Z}\implies\]


Any equivalent response forms are valid, but the one written above is most convenient for analysis.

\par\smallskip\noindent\textbf{Answer:} \textbackslash{}textbf{Answer:} \[\dfrac{\mathcal E(t)}{2}-U_r(0,t)=\dfrac{R-Z}{R+Z}\left(\dfrac{\mathcal E(t)}{2}-U_l(0,t)\right)\]
\end{solcard}

\begin{solcard}{A2}{Find the time dependence of voltage $U_\mathrm{OSC}(t)$ on an oscilloscope.}{1.10}

Let us denote it by $U_n=U_0\left(1-\left(\dfrac{R-Z}{R+Z}\right)^n\right)$ for convenience. At the initial time $U_l(0,0)=0$, $U_r(0,0)=U_1/2$. A wave propagating to the right reaches the oscilloscope in time $L/v$ and is reflected from the edge of the two-wire line. After another time $L/v$, the wave traveling to the left returns to the resistor, and then $U_l(0,2L/v)=U_1/2\implies U_l(0,2L/v)=U_2/2$. Next, the process of re-reflection of the traveling waves is repeated, and after time $2nL/v$ the voltage $U_l(x=0)$ reaches $U_n/2$, and the voltage $U_r(x=0)=U_{n+1}/2$. From here we come to the final result:

\par\smallskip\noindent\textbf{Answer:} \textbackslash{}textbf{Answer:} \[U_\mathrm{OSC}(t)=U_0\left[1-\left(\dfrac{R-Z}{R+Z}\right)^{\left\lfloor\frac{Vt}{2L}+\frac{1}{2}\right\rfloor}\right]\]
\end{solcard}

\begin{solcard}{A3}{By applying a rectangular signal from the generator, remove the dependence of voltage on time. Record the times $t_n$ when the $n$ voltage step occurred, and the values ​​of voltage $U_n$ immediately after.}{0.70}

The voltage step ladder appears on time scales of order $L/c\sim 10^{-8}\ldots10^{-7}~\mathrm{s}$, so it is visible only on time scales of order $\sim100~\mathrm{ns}$. Although the ladder becomes clearer with more cable sections, the number of unlubricated “steps” is still small, so measurements must be taken with a large number of cable sections. A typical view of the ladder is shown in the photo below:


A typical view of the ladder is shown in the photo below:

\begin{figure}[H]
  \centering
  \includegraphics[width=0.72\textwidth]{E2_S_EN_assets/fig1.jpg}
\end{figure}
\end{solcard}

\begin{solcard}{A4}{Repeat the previous step, connecting more cable sections (from 2 to 5 inclusive).}{2.80}

$m=1$ cable section:


$n$ 1 2 3 4 5 6 7 8 $(2n-1)m$ 1 3 5 7 9 11 13 15 $t_n, ~\mathrm{ns}$ 44 95.2 151 205 258 314 365 424 $\Delta t_n,~\mathrm{ns}$ 18 17 19 18 18 19 17 20 $U,~\mathrm{mV}$ -479 -470 -460 -451 -442 -433 -424 -415 $\Delta U,~\mathrm{mV}$ 3 3 3 3 3 3 3 3 $u$ -0.479 -0.47 -0.46 -0.451 -0.442 -0.433 -0.424 -0.415 $\Delta u$ 0.003 0.003 0.003 0.003 0.003 0.003 0.003 0.003


$m=2$ cable section:


$n$ 1 2 3 4 5 6 7 8 $(2n-1)m$ 2 6 10 14 18 22 26 30 $t_n, ~\mathrm{ns}$ 28 130 238 338 448 548 656 756 $\Delta t_n,~\mathrm{ns}$ 35 34 36 33 37 33 36 33 $U,~\mathrm{mV}$ -479 -469 -459 -450 -440 -431 -422 -413 $\Delta U,~\mathrm{mV}$ 3 3 3 3 3 3 3 3 $u$ -0.479 -0.469 -0.459 -0.45 -0.44 -0.431 -0.422 -0.413 $\Delta u$ 0.003 0.003 0.003 0.003 0.003 0.003 0.003 0.003


$m=3$ cable section:


$n$ 1 2 3 4 5 6 7 8 $(2n-1)m$ 3 9 15 21 27 33 39 45 $t_n, ~\mathrm{ns}$ 50 210 364 520 680 834 988 1150 $\Delta t_n,~\mathrm{ns}$ 52 53 51 52 53 51 51 54 $U,~\mathrm{mV}$ -489 -478 -469 -459 -449 -440 -430 -421 $\Delta U,~\mathrm{mV}$ 4 4 3 3 3 3 3 3 $u$ -0.489 -0.478 -0.469 -0.459 -0.449 -0.44 -0.43 -0.421 $\Delta u$ 0.004 0.004 0.003 0.003 0.003 0.003 0.003 0.003


$m=4$ cable section:


$n$ 1 2 3 4 5 6 7 8 $(2n-1)m$ 4 12 20 28 36 44 52 60 $t_n, ~\mathrm{ns}$ 80 284 496 704 906 1120 1320 1530 $\Delta t_n,~\mathrm{ns}$ 69 68 71 69 67 71 67 70 $U,~\mathrm{mV}$ -489 -479 -468 -458 -449 -439 -430 -420 $\Delta U,~\mathrm{mV}$ 3 3 4 3 3 3 3 3 $u$ -0.489 -0.479 -0.468 -0.458 -0.449 -0.439 -0.43 -0.42 $\Delta u$ 0.003 0.003 0.004 0.003 0.003 0.003 0.003 0.003


$m=5$ cable sections:


$n$ 1 2 3 4 5 6 7 8 $(2n-1)m$ 5 15 25 35 45 55 65 75 $t_n, ~\mathrm{ns}$ 104 360 616 876 1120 1390 1660 1900 $\Delta t_n,~\mathrm{ns}$ 84 85 85 87 81 90 90 80 $U,~\mathrm{mV}$ -489 -479 -469 -459 -449 -439 -429 -420 $\Delta U,~\mathrm{mV}$ 3 3 3 3 3 3 3 3 $u$ -0.489 -0.479 -0.469 -0.459 -0.449 -0.439 -0.429 -0.42 $\Delta u$ 0.003 0.003 0.003 0.003 0.003 0.003 0.003 0.003

\end{solcard}

\begin{solcard}{A5}{Neglecting the length of the cables used to connect the generator, linearize the dependences $u(n,m)$ and $t(n,m)$. Plot the linearized dependence $u(n,m)$ and determine the resistance $R$ of the supplied resistor if the characteristic impedance of the cable is $50~ Ом$. Plot a linearized graph of $t(n,m)$ and determine the speed $v$ of wave propagation in the coaxial cable. Give the errors of the results.}{1.00}

For $m$ connected cable sections we will have: \[t(n,m)=(2n-1)mL_0/v,\quad L_0=5~\mathrm{m}.\]Linearization -- $t((2n-1)m)$ with slope $L_0/v$. Since the voltage on the unlubricated part of the ladder is much less than the voltage $U_0$, we can write for it as a first approximation: \[U_n=U_0\left(1-\left(\dfrac{R-Z}{R+Z}\right)^n\right)\approx\dfrac{2nZ}{R}U_0.\]Therefore, a convenient linearization is $u(n)$, and its slope is equal to ${2Z}/{R}$.


The angular coefficient of dependence $u(n)$ is equal to $\dfrac{2Z}{R}=9.6\cdot10^{-3}$, we estimate the error as $\varepsilon_R\sim 2{\Delta u}/(u_{\max}-u_{\min})\approx 9\,\%$, then:


The angular coefficient of dependence $t(m(2n-1))$ is equal to $\dfrac{L_0}{v}=25.4~\mathrm{ns}$, we estimate the error as $\varepsilon_v\sim 2{\Delta t}/(t_{\max}-t_{\min})\approx 9\,\%$, then:

\par\smallskip\noindent\textbf{Answer:} \textbackslash{}textbf{Answer:} \[R=(10.4\pm0.9)~\mathrm{k\Omega}\]
\begin{figure}[H]
  \centering
  \includegraphics[width=0.72\textwidth]{E2_S_EN_assets/fig2.jpg}
\end{figure}
\begin{figure}[H]
  \centering
  \includegraphics[width=0.72\textwidth]{E2_S_EN_assets/fig3.jpg}
\end{figure}
\end{solcard}

\begin{solcard}{B1}{Neglecting the intrinsic capacitance of the cables used to connect the generator and oscilloscope, propose a circuit that allows you to determine the capacitance of a coaxial cable per unit length. Take measurements for different numbers of cable sections and determine its capacity per unit length. Knowing the characteristic impedance of the cable, find its inductance per unit length. Estimate the error of the obtained results.}{1.50}

The most convenient connection option:


We will measure the ratio of the signal amplitudes on the cable (capacitor) and the entire circuit. Since in such a chain: \[\left|\dfrac{U_{RC}}{U_{C}}\right|^2=\left|\dfrac{R+\dfrac{1}{i\omega C}}{\dfrac{1}{i\omega C}}\right|^2=1+4\pi^2R^2L_0^2\mathcal C^2\cdot m^2f^2.\]We use the linearization $\left|\dfrac{U_{RC}}{U_{C}}\right|^2\left(m^2f^2\right)$. The input voltage is equal to $U_{RC}=1~\mathrm{V}$ throughout. Since the error in the resistor resistance value already reaches $5-10\,\%$, we can neglect the oscilloscope error. $m=1$ cable section:


$m=1$ cable section:


$f,~\mathrm{kHz}$ 10 20 30 40 50 60 70 $U_{C}, ~\mathrm{mV}$ 938 810 695 585 504 436 378 $m^2f^2,~10^9~\mathrm{Hz}^2$ 0.1 0.4 0.9 1.6 2.5 3.6 4.9 $(1~\mathrm{V}/U_C)^2$ 1.14 1.52 2.07 2.92 3.94 5.26 7.00


$m=2$ cable section:


$f,~\mathrm{kHz}$ 10 20 30 40 45 50 55 60 $U_{C}, ~\mathrm{mV}$ 821 590 446 344 315 282 260 234 $m^2f^2,~10^9~\mathrm{Hz}^2$ 0.4 1.6 3.6 6.4 8.1 10.0 12.1 14.4 $(1~\mathrm{V}/U_C)^2$ 1.48 2.87 5.03 8.45 10.08 12.57 14.79 18.26


$m=3$ cable section:


$f,~\mathrm{kHz}$ 5 10 15 20 25 30 35 $U_{C}, ~\mathrm{mV}$ 889 723 552 448 368 314 274 $m^2f^2,~10^9~\mathrm{Hz}^2$ 0.23 0.90 2.03 3.60 5.63 8.10 11.03 $(1~\mathrm{V}/U_C)^2$ 1.27 1.91 3.28 4.98 7.38 10.14 13.32


$m=4$ cable section:


$f,~\mathrm{kHz}$ 5 10 13 16 19 22 25 $U_{C}, ~\mathrm{mV}$ 823 600 505 427 368 323 286 $m^2f^2,~10^9~\mathrm{Hz}^2$ 0.40 1.60 2.70 4.10 5.78 7.74 10.00 $(1~\mathrm{V}/U_C)^2$ 1.48 2.78 3.92 5.48 7.38 9.59 12.23


$m=5$ cable sections:


$f,~\mathrm{kHz}$ 3 6 8 10 12 14 16 $U_{C}, ~\mathrm{mV}$ 885 711 596 523 448 388 347 $m^2f^2,~10^9~\mathrm{Hz}^2$ 0.23 0.9 1.6 2.5 3.6 4.9 6.4 $(1~\mathrm{V}/U_C)^2$ 1.28 1.98 2.82 3.66 4.98 6.64 8.31


The slope is $4\pi^2R^2L_0^2\mathcal C^2=1.15\cdot10^{-9}~\mathrm{Hz}^{-2}$, hence $\mathcal C=1.04\cdot10^{-10}~\mathrm{F}/м$. We estimate the error as $\sim10\,\%$, then:


Since $Z=\sqrt{\dfrac{\mathcal L}{\mathcal C}}$, the cable inductance per unit length is $\mathcal L=Z^2\mathcal C=2.6\cdot10^{-7}~\mathrm{H}/м$.

\par\smallskip\noindent\textbf{Answer:} \textbackslash{}textbf{Answer:} \[\mathcal C=\left(1.04\pm0.10\right)\cdot10^{-10}~\mathrm{F}/м\]
\begin{figure}[H]
  \centering
  \includegraphics[width=0.72\textwidth]{E2_S_EN_assets/fig4.jpg}
\end{figure}
\begin{figure}[H]
  \centering
  \includegraphics[width=0.72\textwidth]{E2_S_EN_assets/fig5.jpg}
\end{figure}
\begin{figure}[H]
  \centering
  \includegraphics[width=0.72\textwidth]{E2_S_EN_assets/fig6.jpg}
\end{figure}
\end{solcard}

\begin{solcard}{C1}{Having supplied a rectangular signal from the generator, remove the dependence of voltage on time, namely, find the moments of time $t_n$ when the $n$ stepwise increase in voltage occurred, and the voltage values ​​$U_n$ immediately after that.}{1.00}

$n$ 1 2 3 4 5 6 7 8 $2n-1$ 1 3 5 7 9 11 13 15 $t_n, ~\mathrm{ns}$ 24 134 232 338 436 548 648 750 $\Delta t_n,~\mathrm{ns}$ 110 98 106 98 112 100 102 $U,~\mathrm{mV}$ -488 -468 -447 -427 -408 -389 -371 -353 $\Delta U,~\mathrm{mV}$ 20 21 20 19 19 18 18 $u$ -0.488 -0.468 -0.447 -0.427 -0.408 -0.389 -0.371 -0.353 $\Delta u$ 0.02 0.021 0.02 0.019 0.019 0.018 0.018

\end{solcard}

\begin{solcard}{C2}{Find the speed $v$ of signal propagation in the twisted pair and its characteristic impedance $Z$. Give the errors in the results obtained.}{1.50}

The angular coefficient of dependence $u(n)$ is equal to $\dfrac{2Z}{R}=1.93\cdot10^{-2}$, we estimate the error as $\varepsilon_Z\sim 2{\Delta u}/(u_{\max}-u_{\min})\approx 15\,\%$, then:


The angular coefficient of dependence $t(2n-1)$ is equal to $\dfrac{L_0}{v}=51.8~\mathrm{ns}$, we estimate the error as $\varepsilon_v\sim 2{\Delta t}/(t_{\max}-t_{\min})\approx 9\,\%$, then:

\par\smallskip\noindent\textbf{Answer:} \textbackslash{}textbf{Answer:} \[R=(100\pm20)~\mathrm{\Omega}\]
\begin{figure}[H]
  \centering
  \includegraphics[width=0.72\textwidth]{E2_S_EN_assets/fig7.jpg}
\end{figure}
\begin{figure}[H]
  \centering
  \includegraphics[width=0.72\textwidth]{E2_S_EN_assets/fig8.jpg}
\end{figure}
\end{solcard}

\end{document}